\section{Cadre théorique de la conférence}
Le développement par dialogue déplace une partie du travail de programmation vers la description du problème. Le membre du club dispose d'une connaissance que l'assistant ne peut déduire du code : où un arbre masque l'horizon, quelle image conserve la trace mesurable ou pourquoi deux noms de catalogue désignent la même cible. Les sept thèmes ci-dessous explicitent la manière dont l'atelier mobilise cette connaissance. Les annexes A et B donnent les supports complets et le déroulement oral.

\subsection{Du besoin astronomique à une tâche programmable}
Quatre difficultés motivent le recours à l'IA. La première est le démarrage : une personne sait décrire l'affiche à imprimer sans connaître la syntaxe nécessaire pour la composer. La deuxième est l'intégration : produire une fiche de soirée exige de relier un catalogue, un calcul de coordonnées et un outil de mise en page. Un premier essai utile fait circuler un seul objet entre ces composants. La troisième difficulté est le diagnostic. Le message exact du terminal, la commande lancée et un petit fichier qui reproduit le problème permettent de distinguer une erreur de chemin d'une valeur astronomique incohérente. Enfin, l'automatisation rend envisageables des outils de niche, comme la documentation répétée des poses d'un même instrument. Le gain attendu dépend de la fréquence de la tâche et du temps consacré aux contrôles.

\subsection{Cinq bénéfices à examiner sur le terrain}
Kim et Yegge (2025) regroupent cinq bénéfices sous le sigle FAAFO. La \emph{rapidité} concerne le délai avant un premier résultat testable. L'\emph{ambition} concerne la portée du projet : après avoir traité Messier, on peut envisager un second catalogue. L'\emph{autonomie} tient à la possibilité d'adapter l'outil à son horizon et à ses conventions de classement. Le \emph{plaisir} de retrouver sa propre image sur l'affiche peut soutenir l'effort d'apprentissage. Enfin, la \emph{possibilité de choisir} entre plusieurs réalisations autorise de vrais essais comparatifs : présenter les mêmes données sous forme de galerie ou de rapport imprimable permet au club de juger chaque sortie dans ses conditions d'usage. Ces dimensions proposent des questions d'évaluation ; l'atelier n'en mesure pas les effets sur un groupe de participants.

\subsection{Erreurs silencieuses, régressions et responsabilité}
Une correspondance de catalogue inventée produit parfois une affiche impeccable. Le contrôle scientifique consiste alors à remonter au catalogue de référence. Une régression survient lorsqu'une modification détériore une fonction déjà validée. Dans le récit de GNU Astro Galery, un fichier proposé comme corrigé est passé de 86 à 60~ko. Cette variation a servi d'alerte ; seule la comparaison des versions et des comportements permet d'identifier ce qui a disparu. Les ajustements successifs peuvent aussi créer une dette technique, c'est-à-dire rendre la correction suivante plus coûteuse. Un essai conservé pour chaque anomalie permet de détecter sa réapparition.

Les notes distinguent trois responsabilités : travailler sur des copies pour préserver les acquisitions, garder les clés d'accès hors des fichiers publiés et consigner la provenance du code et des images. Une quatrième responsabilité est cognitive. Si le membre est incapable d'expliquer comment une pose a été sélectionnée, il lui devient difficile de contester une réponse plausible. La relecture par un pair et un cas présentant un résultat connu donnent à l'astronome un point de comparaison indépendant de l'assistant.

\subsection{Logiciel libre, ouverture du code et vie d'un club}
Le logiciel libre met l'accent sur quatre libertés : exécuter le programme pour l'usage choisi (0), étudier et modifier son code source (1), redistribuer des copies (2) et partager des versions modifiées (3). L'approche \emph{open source} met notamment en avant les conditions de redistribution et la possibilité de collaborer sur un code accessible. Les deux traditions se recoupent. Pour le club, la question concrète est de savoir à quelles conditions il pourra adapter l'outil et transmettre sa version au club voisin.

Une licence permissive, comme MIT, facilite la réutilisation sous diverses conditions, sous réserve notamment de conserver les notices requises. Une licence à réciprocité, comme la GPL, impose des conditions de partage lorsqu'on distribue une version couverte. Les crédits des photographies, des catalogues et des bibliothèques doivent être examinés séparément. Eghbal (2020) observe par ailleurs que la répartition du travail varie selon la taille du public et le nombre de contributeurs : elle distingue les outils personnels, les clubs, les stades et les fédérations. Le passage d'un outil personnel à un projet de club devient concret lorsqu'une deuxième personne peut l'installer, reproduire un exemple et proposer une correction.

\subsection{Formuler une consigne, choisir un mode et estimer l'usage}
Une consigne testable précise cinq éléments : le rôle demandé, le contexte, l'objectif, les contraintes et la forme du résultat. Pour l'affiche, cela pourrait donner : « Aide un membre débutant à classer trois JPEG de son Seestar sous Windows. Associe M31 et NGC~224 à une seule cible. Conserve les originaux. Rends un tableau indiquant, pour chaque fichier, le nom lu et la cible retenue. » L'observateur peut valider ce tableau avant de commander la mise en page des 110 cases.

Les notes proposent cinq façons de lever une ambiguïté : fournir quelques exemples réels, demander un plan observable, inviter l'assistant à poser une question, lui faire relire sa proposition selon un critère précis et reformuler la consigne avant exécution. Une demande qui mêle lecture FITS, météo, carte et export PDF multiplie les causes possibles d'échec. Les tester successivement permet de localiser une erreur. Quand le dialogue tourne en rond, une nouvelle demande peut repartir du fichier fautif, du message exact et du dernier résultat validé.

L'atelier illustre trois usages : explorer les besoins dans Chat, préparer un document dans Work, puis modifier le dépôt et exécuter ses vérifications dans Codex. L'unité de texte traitée par un modèle est le \emph{jeton} (\emph{token}). La fenêtre de contexte borne la quantité d'information que le modèle peut prendre en compte au cours d'un appel. Une tâche sur plusieurs fichiers peut mobiliser davantage de jetons qu'une question courte, notamment si elle nécessite plusieurs lectures, essais et réponses.

Pour une facturation à l'usage, on multiplie séparément les volumes d'entrée ordinaire, d'entrée en cache et de sortie par leurs tarifs respectifs, puis on additionne ces produits. Lorsque le tarif est donné par million de jetons, chaque volume est divisé par un million avant la multiplication. Une formule d'abonnement peut, elle, appliquer des limites d'usage. Le volume dépend de la tâche, du contexte transmis et du modèle. Les exemples de prix et de quotas que le présentateur commentait sont conservés dans l'annexe B avec leur contexte temporel.

\subsection{Anatomie d'un projet réutilisable}
À la racine, le fichier \texttt{README.md} explique l'installation et le premier essai ; la licence expose les conditions de réutilisation. Le dossier \texttt{src/} contient le code, \texttt{tests/} les essais reproductibles et \texttt{docs/} les guides. Le fichier \texttt{.gitignore} écarte du suivi les produits temporaires et les fichiers locaux désignés. Un petit jeu d'exemples dont l'issue est connue aide davantage à vérifier le programme qu'une archive volumineuse dépourvue de résultat attendu.

Les fichiers Markdown répondent à des lecteurs différents. Le journal \texttt{CHANGELOG.md} décrit les changements déjà livrés. La feuille de route \texttt{ROADMAP.md} présente les travaux envisagés. Les consignes \texttt{CONTRIBUTING.md} expliquent à un membre comment signaler un problème ou soumettre une correction. Un code de conduite précise les attentes dans les échanges. Enfin, \texttt{AGENTS.md} conserve les règles de travail de l'assistant, par exemple la préservation des FITS originaux. Ces documents transmettent les décisions prises pendant le dialogue à la personne qui reprendra le projet.

\subsection{Git, GitHub et versions vérifiables}
Git enregistre l'histoire du code sur l'ordinateur du membre. GitHub héberge un dépôt partagé où le club peut examiner les changements et en discuter. Un \emph{commit} mémorise les modifications sélectionnées avec un message. Dans l'exemple suivant, le membre inspecte son travail, sélectionne le guide modifié et enregistre cette seule étape.
\begin{verbatim}
git status
git add README.md
git commit -m "Documenter le premier essai"
\end{verbatim}
De petits commits facilitent la recherche de la modification qui a introduit une erreur. Une branche permet, par exemple, d'essayer une autre règle de tri et de comparer les affiches obtenues. Une fois les essais concluants, la fusion intègre la modification à la branche principale. Le partage du dépôt permet à un second observateur de reprendre la même révision.

Le versionnement sémantique distingue changement majeur, ajout mineur compatible et correctif compatible. À partir de 1.2.0, une correction peut donner 1.2.1, une nouvelle fonction compatible 1.3.0 et une modification incompatible de l'interface publique 2.0.0. Une série 0.x indique un développement initial dont l'interface peut encore évoluer. Inscrire la version exacte de GNU Astrolist dans un carnet d'observation relie le résultat au comportement du logiciel employé ce soir-là.
\nocite{ref02,ref03,ref05,ref09,ref10,ref11,ref12,ref13,ref14,ref15,ref16,ref17}
